\documentclass[10pt,a4paper]{amsart}
\usepackage[margin=19mm]{geometry}
\usepackage{amsmath,amssymb,mathtools}
\usepackage[scr=rsfs,cal=euler]{mathalfa}
\usepackage{xfrac}
\usepackage[unicode,hidelinks,bookmarks=false]{hyperref}
\newcommand{\ie}{i.e.\ }
\newcommand{\cf}{cf.\ }
\newcommand{\resp}{respectively\ }
\newcommand{\Subsection}{\subsection}
\providecommand{\nobreakdash}{\nobreak\hbox{-}}
\setlength{\emergencystretch}{2em}
\multlinegap0pt
\usepackage{amssymb}
\usepackage{mathtools}
\usepackage[shortlabels]{enumitem}
\newtheorem{lemma}{Lemma}
\newtheorem{proposition}{Proposition}
\newcommand{\Bur}{\overline B}
\newcommand{\Th}{\operatorname{Th}}
\newcommand{\dd}{\doteq}
\title{Fox's Trapezoidal Conjecture for Four-Strand Turk's Head Knots and Links}
\author{Suman Saurabh}
\date{Source: arXiv:2606.15256v2 (CC BY 4.0)}
\begin{document}
\maketitle

\noindent\emph{Source and licence.} Fox's Trapezoidal Conjecture for Four-Strand Turk's Head Knots and Links. Version arXiv:2606.15256v2 (CC BY 4.0), distributed by arXiv under the Creative Commons Attribution 4.0 International licence, http://creativecommons.org/licenses/by/4.0/, as recorded in the arXiv licence metadata for this exact version and shown on the abstract page https://arxiv.org/abs/2606.15256v2. Full licence notice: \texttt{licenses/arxiv-2606.15256-CC-BY-4.0.txt}.\par\smallskip

\noindent\textit{Editorial scope.} Counted unit: the article's proposition prop:uniform-factorization (source0.tex lines 236--245). The article's complete original proof of it is delivered with it (source0.tex lines 247--301, one \texttt{proof} environment of 55 lines). No mathematics is written, repaired or re-derived: the statement and the proof are the article's own lines, in the article's own notation, and no retained line is substituted or re-flowed.  Retained uncounted as the source-local context the proof needs: lines 157--161; lines 210--216.  Nothing was omitted from any retained span, so the package carries no omission record.  No retained line contains a citation, so the wrapper carries no bibliography and no bibliography entry.  No retained line contains a figure environment, an image-inclusion command or drawing code, so no image is delivered and the package declares no asset dependency.  Editorial formatting only, in addition: an AMS \texttt{amsart} preamble that restates the article's own declarations and macros used by the retained lines, namely the environments \texttt{lemma}, \texttt{proposition} on separate counters, together with the standard packages those lines need.  The retained environments print in this document as Lemma 1, Proposition 1 in order of appearance, whereas the article prints the same items with the numbering of its own full document.  The complete native source of the article as distributed in the arXiv e-print for this exact version is bundled unchanged as \texttt{sources/202845/source0.tex}.
\begin{equation}\label{eq:burau-formula-general}
        \Delta_{\widehat{\gamma}}(t)
        \dd
        \frac{1-t}{1-t^p}\det\left(I-\Bur(\gamma)\right).
\end{equation}

\begin{lemma}\label{lem:det-I-xM}
For an auxiliary variable \(x\),
\begin{equation}\label{eq:det-I-xM}
        \det(I-xM(t))
        =\frac{(1+tx)\bigl(tx^2+(t-1)^2x+t\bigr)}{t}.
\end{equation}
\end{lemma}

\begin{proposition}[Uniform spectral factorization]\label{prop:uniform-factorization}
For every \(q\ge1\),
\begin{equation}\label{eq:uniform-factorization}
        \Delta_{\Th(4,q)}(-z)
        \dd
        (1+z+\cdots+z^{q-1})
        \prod_{\substack{\omega^q=1\\ \omega\ne1}}
        \left(z^2+(2-\omega-\omega^{-1})z+1\right).
\end{equation}
\end{proposition}

\begin{proof}
By \eqref{eq:burau-formula-general},
\[
        \Delta_{\widehat{\beta^q}}(t)
        \dd
        \frac{1-t}{1-t^4}\det(I-M(t)^q).
\]
For any square matrix \(M\),
\[
        \det(I-M^q)=\prod_{\omega^q=1}\det(I-\omega M).
\]
Indeed, the scalar identity \(1-X^q=\prod_{\omega^q=1}(1-\omega X)\) gives the corresponding matrix identity because all factors are polynomials in \(M\); taking determinants gives the displayed formula.

By Lemma~\ref{lem:det-I-xM},
\[
        \det(I-\omega M(t))
        =\frac{(1+t\omega)\bigl(t\omega^2+(t-1)^2\omega+t\bigr)}{t}.
\]
Now set \(t=-z\).  Then
\[
\begin{aligned}
        t\omega^2+(t-1)^2\omega+t
        &=-z\omega^2+(1+z)^2\omega-z  \\
        &=\omega\left(z^2+(2-\omega-\omega^{-1})z+1\right).
\end{aligned}
\]
Therefore
\begin{equation}\label{eq:single-root-factor}
        \det(I-\omega M(-z))
        =-\frac{\omega}{z}(1-z\omega)
        \left(z^2+(2-\omega-\omega^{-1})z+1\right).
\end{equation}
The prefactor \(-\omega/z\), when multiplied over all \(q\)-th roots of unity, contributes only a Laurent unit in \(z\).  Hence
\[
\begin{aligned}
\Delta_{\widehat{\beta^q}}(-z)
&\dd
\frac{1+z}{1-z^4}
\prod_{\omega^q=1}(1-z\omega)
\prod_{\omega^q=1}
\left(z^2+(2-\omega-\omega^{-1})z+1\right)  \\
&=
\frac{1+z}{1-z^4}(1-z^q)(z^2+1)
\prod_{\substack{\omega^q=1\\ \omega\ne1}}
\left(z^2+(2-\omega-\omega^{-1})z+1\right).
\end{aligned}
\]
Since \(1-z^4=(1-z)(1+z)(1+z^2)\), this becomes
\[
        \frac{1-z^q}{1-z}
        \prod_{\substack{\omega^q=1\\ \omega\ne1}}
        \left(z^2+(2-\omega-\omega^{-1})z+1\right),
\]
which is \eqref{eq:uniform-factorization}.
\end{proof}

\end{document}
